% CAPÍTULO 6 -- CONCLUSÃO (VERSÃO CONCISA TCC UESC)
% ==========================================================================
\chapter{Conclusão}
\label{cap:conclusao}

\section{Síntese das realizações do trabalho}
\label{sec:conc-sintese}

O presente trabalho de conclusão de curso demonstrou a viabilidade técnica, a superioridade elastodinâmica e a qualificação aeroespacial de um chassi primário para nanossatélite CubeSat 1U integrado por metamateriais celulares auxéticos em liga de alumínio AlSi10Mg aliviada termicamente (\SI{300}{\celsius}/\SI{2}{\hour}), fabricado por Manufatura Aditiva via Fusão em Leito de Pó por Laser (L-PBF).

A metodologia proposta integrou com sucesso uma esteira computacional de ponta a ponta totalmente aberta em ambiente Linux, unificando a modelagem geométrica analítica B-Rep no OpenCASCADE (CadQuery / build123d), o gerador de malhas tetraédricas de alta ordem Gmsh (Tet10 com até \num{342800} nós), solvers de elementos finitos abertos (Code\_Aster e CalculiX), metamodelagem neural profunda informada por física (\textit{Physics-Guided ResNet}), otimização multiobjetivo pelo algoritmo NSGA-II e tomada de decisão multicritério por TOPSIS.

A configuração de voo ótima identificada (\textbf{Projeto \#28}), caracterizada por ângulo de inclinação $\theta^* = \SI{65.75}{\degree}$ (com deflexão cartesiana $\theta_{\text{def}}^* = \SI{-21.40}{\degree}$ e reentrância $\phi^* = \SI{24.25}{\degree}$), espessura de costela $t^* = \SI{0.613}{\milli\metre}$ ($t_{\text{mech,eff}} \approx \SI{0.563}{\milli\metre}$), comprimento oblíquo $l^* = \SI{10.42}{\milli\metre}$ e altura vertical $h^* = \SI{14.85}{\milli\metre}$, com Poisson micromecânico 2D $\nu_{yx}^* = \num{-12.63}$ e Poisson estrutural 3D homogeneizado $\nu_{\text{eff}} = \num{-0.428}$, satisfez rigorosamente todos os portões de DfAM e as exigências normativas da NASA GSFC-STD-7000A, NASA-STD-5020B e NASA-HDBK-7005.

\section{Validação quantitativa das hipóteses de trabalho}
\label{sec:conc-hipoteses}

Os resultados computacionais comprovaram plenamente as três hipóteses de trabalho formuladas na introdução:

\begin{enumerate}
  \item \textbf{Validação da Hipótese 1 ($H_1$ -- Alívio mássico com preservação de rigidez):}
  \begin{itemize}
    \item A massa da estrutura primária foi reduzida de \SI{0.308}{\kilo\gram} (chassi monolítico usinado convencional) para \SI{0.11554}{\kilo\gram} (arredondado para \SI{0.116}{\kilo\gram}), alcançando um **alívio de massa de \SI{62.49}{\percent}** (nominalmente \SI{192.46}{\gram} de economia estrutural), superando amplamente a meta de \SI{40.0}{\percent};
    \item A primeira frequência natural fundamental sob engastamento rígido atingiu $f_1 = \SI{579.6}{\hertz}$ (\SI{+13.20}{\percent} superior aos \SI{512.0}{\hertz} do monolítico) e $f_1 = \SI{304.1}{\hertz}$ sob suporte elástico no dispensador P-POD (\SI{+13.30}{\percent} superior aos \SI{268.4}{\hertz} do monolítico), garantindo Margem de Segurança de rigidez $MS_{f_1} = +2{,}04$ em relação ao piso normativo aeroespacial ($f_1 \ge \SI{100.0}{\hertz}$).
  \end{itemize}

  \item \textbf{Validação da Hipótese 2 ($H_2$ -- Isolamento dinâmico passivo da carga útil):}
  \begin{itemize}
    \item Sob a excitação estocástica normativa de \SI{14.1}{\Grms} da NASA GEVS, operando no regime sub-comprimento de onda ($a/\lambda \le \num{0.006} \ll 1$), a aceleração eficaz na interface da carga útil foi reduzida para apenas $\Grmspay = \SI{2.82}{\Grms}$, correspondendo a uma transmissibilidade dinâmica de $T = \num{0.1997}$;
    \item Esse resultado comprova uma **atenuação elastodinâmica passiva de \SI{76.50}{\percent} (\SI{12.6}{\decibel})**, viabilizada por descasamento severo de impedância acústica ($Z_{\text{solid}} = \SI{13.47}{\mega\pascal\second\per\metre}$ vs $Z_{\text{eff}} = \SI{1.39}{\mega\pascal\second\per\metre}$, reflexão $> \SI{66.1}{\percent}$), antirressonâncias mecânicas passa-baixa e desacoplamento de modos de tombamento (\textit{rocking}) via Poisson negativo, sem necessidade de amortecedores poliméricos externos sujeitos a desgaseificação (\textit{outgassing}).
  \end{itemize}

  \item \textbf{Validação da Hipótese 3 ($H_3$ -- Eficiência de metamodelagem e confiabilidade à fadiga):}
  \begin{itemize}
    \item A \textit{Physics-Guided ResNet} atingiu coeficientes de determinação $\overline{R^2} = \num{0.9972} \pm \num{0.0016}$ e erro $\text{NRMSE} = \SI{1.27 \pm 0.34}{\percent}$ sob validação cruzada estratificada em 10 dobras ($10$-fold CV), operando sobre vetor de saída desacoplado em $\mathbb{R}^3$ e eliminando redundâncias lineares, proporcionando um \textit{speedup} instantâneo de inferência superior a \num{200000} vezes frente ao solver 3D direto e um \textit{speedup} amortizado de ponta a ponta de $\approx \num{40.0}\times$ (avaliando \num{10000} topologias em \SI{4.82}{\second});
    \item A tensão equivalente de pico a $3\sigma$ estabilizou em \SI{166.91}{\mega\pascal}, estritamente abaixo da tensão admissível de escoamento ($\sigma_{\text{adm, yield}} = \SI{184.0}{\mega\pascal}$), assegurando Margens de Segurança positivas: $MS_{\text{yield}} = \num{+0.102} > 0$ e $MS_{\text{ult}} = \num{+0.455} > 0$;
    \item O índice de dano cumulativo por fadiga de Steinberg, avaliado no entalhe crítico de máxima concentração de tensões da costela auxética, alcançou $D = \num{0.0458} \ll \num{0.25}$, garantindo margem de vida de $5{,}46\times$ frente ao ambiente vibratório de lançamento (mantendo-se em $D = \num{0.082} \ll \num{0.25}$ mesmo sob \textit{knockdown} de \SI{10}{\percent} de espessura por rugosidade L-PBF);
    \item O estudo de convergência de malha pelo GCI de Roache (ASME V\&V 20) comprovou incerteza de malha de apenas $GCI = \SI{0.23}{\percent}$ para $f_1$, o Método do Pico de Tensão (PSM/NSIF) atestou integridade subcrítica contra trincas ($\Delta K_1^V = \SI{1.12}{\mega\pascal\sqrt{\metre}} < \Delta K_{\text{th}} = \SI{2.50}{\mega\pascal\sqrt{\metre}}$) e a tolerância a porosidades estocásticas de L-PBF via mecânica da fratura de Murakami ($\Delta K_{\text{defect}} = \SI{1.28}{\mega\pascal\sqrt{\metre}} < \Delta K_{\text{th}}$, $a_{\text{th}} \approx \SI{169.0}{\micro\metre} \gg \SI{50.0}{\micro\metre}$) chancelou a integridade elástica sob rota térmica dual ($\SI{300}{\celsius}$ e $\SI{250}{\celsius}$);
  \end{itemize}
\end{enumerate}

\section{Contribuições científicas e tecnológicas}
\label{sec:conc-contribuicoes}

As principais contribuições deste trabalho para o estado da arte e para o desenvolvimento aeroespacial compreendem:
\begin{alineas}
  \item **Demonstração pioneira de sinergia entre metamateriais e DfAM aeroespacial:** prova de conceito numérica rigorosa de que painéis auxéticos reentrantes podem ser plenamente qualificados para voo espacial em L-PBF sem suportes internos;
  \item **Desenvolvimento de esteira computacional de alta fidelidade 100\% aberta:** estabelecimento de um fluxo CAE completo em Linux (OpenCASCADE, Gmsh, Code\_Aster, CalculiX), dispensando licenças comerciais de software e viabilizando pesquisas de ponta em instituições públicas com orçamento restrito;
  \item **Inovação em metamodelagem informada por física:** demonstração prática de que a introdução de funções de perda baseadas nas leis de escala de Gibson-Ashby e monotonicidade de rigidez elimina predições espúrias em redes neurais profundas aplicadas a sólidos celulares;
  \item **Maturidade tecnológica comprovada:** consolidação do patamar de prontidão \textbf{Digital TRL 4} (NASA SP-6105) para o arcabouço computacional analítico e preditivo validado sob o perfil NASA GEVS, e \textbf{TRL 3 Avançado de Hardware} (ISO 16290:2013) para os modelos físicos de interface.
\end{alineas}

\section{Trabalhos futuros}
\label{sec:conc-trabalhos-futuros}

Como desdobramentos naturais desta pesquisa rumo à qualificação de voo orbital (TRL 5 a 7), recomendam-se as seguintes linhas de investigação:
\begin{enumerate}
  \item **Fabricação física e ensaios experimentais de bancada:** imprimir o Projeto \#28 em máquina industrial L-PBF em liga AlSi10Mg e submeter a estrutura a ensaios dinâmicos em mesa vibratória (\textit{shaker}) de \SI{20}{\kilo\newton} com perfil NASA GEVS, superando a similaridade incompleta inerente aos modelos preliminares de empacotamento em PLA;
  \item **Análise termoelástica e câmara de termovácuo (TVAC):** modelar o comportamento da matriz auxética sob ciclos térmicos severos em órbita baixa (\SIrange{-40}{+85}{\celsius}) em termovácuo ($p \le \SI{e-6}{\torr}$), avaliando gradientes de dilatação térmica e condutividade aparente;
  \item **Gradação funcional tridimensional (FG-Lattice):** investigar células auxéticas com variação espacial contínua de densidade e espessura ao longo da altura do chassi, direcionando a absorção de energia para os pontos de fixação dos subsistemas mais críticos;
  \item **Ensaios de separação pirotécnica e choque no P-POD:** caracterizar experimentalmente a capacidade de amortecimento de choques de ejeção do dispensador P-POD (\textit{pyroshock attenuation}) através de testes instrumentados com acelerômetros piezoelétricos de alta frequência.
\end{enumerate}
